% !TEX program = pdflatex
% =====================================================================
% ICCSCI 2026 — Procedia Computer Science
% SLR Paper: Machine Learning for Village Fund Anomaly Detection:
%   A Systematic Literature Review of the Operationalization Chasm
%   in Indonesian Anti-Corruption Information Systems
%
% Template: elsarticle + ecrc (Elsevier Procedia CRC)
% Citation style: Vancouver numbered [#]
% Double-blind submission — all author identities removed
% =====================================================================

\documentclass[3p,times,procedia]{elsarticle}
\flushbottom

%% ecrc package — required for Procedia CRC layout
\PassOptionsToPackage{procedia}{ecrc}
\usepackage{ecrc}
\usepackage[bookmarks=false]{hyperref}
    \hypersetup{colorlinks,
      linkcolor=blue,
      citecolor=blue,
      urlcolor=blue}

%% Core packages
\usepackage{graphicx}
\usepackage{stfloats}
\usepackage{amsmath}
\usepackage{amssymb}
\usepackage{booktabs}
\usepackage{array}
\usepackage{tabularx}

\graphicspath{{../../analysis/bibliometric/}}

%% Journal metadata
\volume{00}
\firstpage{1}
\journalname{Procedia Computer Science}
\runauth{[BLINDED FOR REVIEW]}
\jid{procs}

% =====================================================================
\begin{document}
\begin{frontmatter}

\dochead{The 13th International Conference on Computer Science and Computational Intelligence (ICCSCI 2026)}

\title{Machine Learning for Village Fund Anomaly Detection: A Systematic Literature Review of the Operationalization Chasm in Indonesian Anti-Corruption Information Systems}

\author[a]{Yodihartomo, Farrell}
\author[b]{Faiq, Humam}
\author[a]{Wicaksono, Pandu}
\author[a]{Yen, Charlie\corref{cor1}}
\address[a]{Bina Nusantara University}
\address[b]{Komisi Pemberantasan Korupsi}
\cortext[cor1]{Corresponding author.}

\begin{abstract}
Village fund corruption in Indonesia represents a persistent governance failure: between
2015 and 2023, the Corruption Eradication Commission (KPK) documented over 600 Dana Desa
cases totalling more than IDR 433 billion \cite{method_kpk2023}. Despite accelerating
advances in machine learning (ML)-based financial fraud detection
\cite{p001,p003,p015,p037}, no existing study has applied these methods to the specific
institutional and data conditions of Indonesian village fund governance. This systematic
literature review (SLR) maps and critically synthesises the existing evidence base to
identify what ML methods exist, what feature engineering frameworks have been validated,
and what structural gaps preclude direct application to Dana Desa anomaly detection.
Following PRISMA 2020 guidelines \cite{method_prisma2020}, the review screened 1,001
records and included 45 papers (2018--2026). Thematic synthesis
\cite{method_thomas2008}, bibliometric cluster analysis \cite{method_donthu2021}, and a
Design Science Research (DSR) framework matrix \cite{method_hevner2004} reveal four
analytical themes: an Operationalization Chasm (no study bridges ML detection with village
governance), a Scalability Illusion (performance claims conditioned on inapplicable
assumptions), an Absence of IS Theory (62\% of technical papers contain no IS-theoretical
grounding), and a Ground Truth Paradox (circular synthetic-data validation). Three critical
research voids directly motivate a primary empirical study applying an unsupervised
ensemble to real Jambi Province Dana Desa disbursement data (2023--2025).
\end{abstract}

\begin{keyword}
village fund corruption \sep machine learning \sep anomaly detection \sep
systematic literature review \sep information systems governance
\end{keyword}

\end{frontmatter}
\email{[email-blinded-for-review@institution.edu]}

\vspace*{-6pt}

% =====================================================================
\section{Introduction}
\label{sec:intro}
% =====================================================================

Village fund governance in Indonesia operates at a scale that makes it both critical and
vulnerable. The Village Law (UU No.~6/2014) channels fiscal allocations to over 74,000
village governments annually --- a programme exceeding IDR~70 trillion per fiscal year.
This fiscal decentralisation model places budget execution authority with village officials
who frequently operate outside the supervisory reach of formal audit institutions. The KPK
recorded 601 Dana Desa corruption cases between 2015 and 2023 \cite{method_kpk2023},
implicating village officials in embezzlement, fictitious project expenditure, and budget
absorption manipulation. The financial exposure is systemic, not episodic: recurrence
across consecutive fiscal years in the same villages suggests structural principal-agent
failure rather than isolated opportunistic behaviour \cite{method_jensen1976}.

Current governance responses rely primarily on ex-post audit cycles, which detect fraud
after disbursement has occurred. The capacity gap is structural: Inspektorat Daerah
(APIP) units audit a small fraction of villages annually, and BPK sampling is confined
to aggregate financial statements rather than activity-level expenditure records. This
leaves the operationally critical data layer --- Siskeudes activity transactions ---
largely unscrutinised in real time. ML-based financial fraud detection has accumulated
substantial evidence over the past decade, with Isolation Forest, Local Outlier Factor
(LOF), graph neural networks (GNNs), and autoencoder-based methods reporting AUC-ROC
values exceeding 0.90 on benchmark datasets \cite{p003,p007,p015,p022,p026,p041,p043}.
Yet this technical literature addresses fundamentally different problem conditions:
centralised commercial systems with labelled transaction data. Whether these methods adapt
to village governance conditions --- fragmented records, unlabelled data, sub-national
institutional variation --- is an open question no existing study addresses.

This SLR pursues three research questions:
\begin{itemize}
\item \textbf{RQ1}: What ML-based anomaly detection methods exist in the financial domain
and what performance characteristics do they demonstrate across data and institutional
contexts?
\item \textbf{RQ2}: What feature engineering frameworks operationalize corruption as
machine-detectable computational signals, and how transferable are they to public-sector
governance?
\item \textbf{RQ3}: What structural gaps prevent application of ML detection methods to
decentralised village fund governance, and what design requirements must a primary study
address?
\end{itemize}

This review adopts an IS perspective that frames ML anomaly detection as an IS governance
intervention --- a designed artefact embedded within institutional processes to support
human decision-making \cite{method_hevner2004}. This framing carries a critical
implication: complete artefact evaluation must assess not only predictive accuracy but
also institutional deployability, interpretability, and theoretical alignment with
principal-agent corruption dynamics \cite{method_jensen1976}. The contribution of this
SLR is therefore not merely a gap catalogue --- it is a structured evidence base that
specifies what a viable primary study must design, how it must validate, and why the
two-tradition divide constitutes a theoretically tractable --- not merely empirical ---
problem.

% =====================================================================
\section{Related Work}
\label{sec:related}
% =====================================================================

\subsection{ML-Based Financial Fraud Detection}

Contemporary reviews document three dominant method families: ensemble tree-based
classifiers, deep learning architectures (LSTM, Autoencoder, Transformer), and
graph-based methods \cite{p001,p003,p004,p022,p037}. Unsupervised methods --- Isolation
Forest \cite{p015}, LOF \cite{p041}, and One-Class SVM --- have gained traction where
ground truth labels are unavailable \cite{p003,p015,p041}. The practical implication is
not merely methodological variety: it is that label availability determines which entire
method families are viable. In the Dana Desa context, where confirmed fraud convictions
exist in isolated KPK case files rather than structured database labels, supervised methods
that constitute the majority of high-AUC-ROC reports \cite{p003,p007,p022,p026,p039}
are operationally inaccessible without prior unsupervised or semi-supervised labelling.

Existing SLRs share a critical boundary condition they rarely make explicit: they evaluate
performance on private-sector datasets --- credit card transactions, insurance claims, AML
networks --- and extrapolate generalizability conclusions without establishing whether the
institutional and data assumptions hold in public governance contexts
\cite{p001,p022,p037}. This extrapolation constitutes the epistemological core of the
Scalability Illusion identified as AT2 in this synthesis. Procurement fraud detection
\cite{p008,p043} represents the closest analogue to village fund expenditure anomalies,
but even these studies address centralised national procurement systems, not decentralised
village-level disbursement with administrative capacity constraints.

\subsection{Village Fund Governance and Corruption Patterns}

Dana Desa governance literature documents a consistent corruption typology grounded in
agency theory \cite{method_jensen1976}, the fraud triangle, and institutional control
frameworks: budget absorption manipulation, fictitious procurement, and administrative
fraud \cite{method_kpk2023}. The governance tradition's contribution extends beyond
cataloguing modus operandi: it establishes the institutional conditions --- weak oversight,
rotating village heads, limited audit capacity --- that transform latent corruption
opportunity into observed behaviour. Digital governance studies \cite{p013,p016,p039,p044}
provide IS framing for how technology intersects with these institutional dynamics but
contain no computational detection methods. This omission reveals a division of
intellectual labour that this SLR characterises as structurally problematic: the
community with the richest understanding of corruption mechanisms (governance) has not
attempted to operationalize that understanding as computable features; the community with
the strongest detection algorithms (ML) has not sought to apply them to governance
contexts where the fraud mechanisms differ fundamentally from commercial transaction fraud.

\subsection{Theoretical Framework: DSR and IS Success}

This review organises synthesis through the three-cycle DSR model
\cite{method_hevner2004}: the relevance cycle (village corruption as governance
requirement), the design cycle (ML artefact calibrated to village data), and the rigor
cycle (validated prior knowledge). The model carries a normative implication for this
review: an ML detection tool that satisfies RIGOR (technically sound algorithm) and
RELEVANCE (responsive to corruption typology) but fails the DESIGN cycle (never evaluated
on real village data) cannot be claimed to contribute IS knowledge. The DSR framework
matrix in Section~\ref{sec:results} operationalises this standard and reveals that the
entire literature fails the DESIGN$\times$village criterion.

The DeLone and McLean IS Success Model \cite{method_delone2003} supplements DSR evaluation
by providing artefact assessment dimensions beyond AUC-ROC: information quality
(anomaly scores must be interpretable), system quality (must integrate with Siskeudes
workflows), and net benefits (must reduce APIP inspection workload). Agency Theory
\cite{method_jensen1976} provides the principal-agent corruption mechanism grounding,
establishing that the principal's information asymmetry --- not simply the agent's
opportunism --- is the structural variable a detection system must reduce.

Prior SLRs \cite{p001,p022,p037} cover ML fraud detection broadly but none address
sub-national public governance, feature engineering for village fund corruption typology,
or DSR-based deployability assessment.

% =====================================================================
\section{Methodology}
\label{sec:method}
% =====================================================================

This review follows PRISMA 2020 \cite{method_prisma2020} and adopts Thomas and Harden's
thematic synthesis \cite{method_thomas2008}. Complementary evidence derives from
bibliometric analysis \cite{method_donthu2021} and a DSR framework matrix
\cite{method_hevner2004}. SLR reporting transparency follows Templier and Par\'e (2018)
\cite{method_templier2018}.

\subsection{Search Strategy and Eligibility}

Five databases were searched: OpenAlex (API-automated), Scopus, IEEE Xplore, Web of
Science, and Semantic Scholar. Two complementary string sets were applied: a primary
set (ML methods + public finance fraud, feeding RQ1/RQ2) and a supplementary set (village
fund governance, feeding RQ3). Initial retrieval yielded 1,001 unique records.

Inclusion criteria required: ML or AI methods for financial anomaly/fraud detection
(IC-01); public sector governance or corruption studies (IC-02); IS-theoretical frameworks
(IC-03); publication year 2010--2026 (IC-04); peer-reviewed format (IC-05); and English
or Indonesian language (IC-06). A \textit{domain-override protocol} retained 12 borderline
papers (quality score 4.0--4.99) with direct Dana Desa relevance, documented before
analysis began.

\subsection{Quality Assessment and Selection}

Each paper received a composite quality score (0--10) across five dimensions:
methodological rigour, theoretical grounding, contextual relevance, reproducibility, and
evidence strength. Journal SCImago quartile enrichment was applied to 194 records.
Inclusion threshold: $\geq 5.0$; domain-override threshold: $\geq 4.0$ with confirmed
Dana Desa relevance.

Stage 1 inter-coder reliability reached $\kappa = 0.307$ initially; after adjudication,
$\kappa = 1.000$. Final included corpus: \textbf{N\,=\,45}
\cite{p001,p002,p003,p004,p005,p006,p007,p008,p009,p010,p011,p012,p013,p014,p015,
p016,p017,p018,p019,p020,p021,p022,p023,p024,p025,p026,p027,p028,p029,p030,
p031,p032,p033,p034,p035,p036,p037,p038,p039,p040,p041,p042,p043,p044,p045}.

\subsection{PRISMA Flow}

\begin{table}[htb]
\caption{PRISMA 2020 Study Selection Flow}
\label{tab:prisma}
\begin{tabular*}{\hsize}{@{\extracolsep{\fill}}lr@{}}
\toprule
\textbf{Stage} & \textbf{Records} \\
\colrule
Identification (5 databases) & 1,001 \\
Title/abstract screening (Stage 1) & 1,001 \\
Excluded at Stage 1 & $-888$ \\
Full-text quality assessment & 113 \\
Excluded at Stage 2 & $-68$ \\
CONSENSUS INCLUDE ($\geq$5.0) & 33 \\
DOMAIN\_OVERRIDE (4.0--4.99) & +12 \\
\textbf{Final included corpus} & \textbf{45} \\
\botrule
\end{tabular*}
\end{table}

\subsection{Synthesis Methods}

\textbf{Primary}: Thematic synthesis \cite{method_thomas2008} — open coding (613 code
instances across 45 papers) $\to$ 10 descriptive themes (DT1--DT10) $\to$ 4 analytical
themes (AT1--AT4), with 121 inter-paper relations classified following Britten et al.
(2002) as converging, extending, contradicting, silencing, or bridging.

\textbf{Structural}: DSR framework matrix \cite{method_hevner2004} — papers mapped to
three cycles (DESIGN, RELEVANCE, RIGOR) $\times$ six context levels (village,
sub-national, national, cross-national developing, cross-national developed, unspecified).
Empty cells constitute structural gaps.

\textbf{Complementary}: Bibliometric analysis \cite{method_donthu2021} — temporal trend,
journal distribution, keyword co-occurrence clustering.

\subsection{Sensitivity Analysis}

Three quality thresholds were tested: T1 ($\geq$4.0, N=45), T2 ($\geq$4.5, N=23),
T3 ($\geq$5.0, N=14) \cite{method_higgins2011}.

% =====================================================================
\section{Results}
\label{sec:results}
% =====================================================================

\subsection{Bibliometric Characteristics}

The 45 papers span 2018--2026: 1 paper (2021), 6 (2022), 9 (2023), 12 (2024), 15 (2025),
2 (2026) --- 98\% published 2022--2026, confirming accelerating research output consistent
with post-pandemic intensification of public financial integrity concerns.
IEEE Access (Q1, 9 papers) and Applied Sciences (Q2) dominate the ML detection strand.
The governance strand concentrates in regional IS and public administration journals,
many unranked in Scopus --- a structural asymmetry that partially explains the absence
of cross-citing between traditions. Figure~\ref{fig:trend} shows the temporal distribution;
Figure~\ref{fig:cluster} shows the bibliometric cluster separation.

\begin{figure}[htb]
\centering
\includegraphics[width=0.48\textwidth]{publication_trend.png}
\caption{Publication trend 2021--2026 with cumulative overlay (N\,=\,45).}
\label{fig:trend}
\end{figure}

\begin{figure}[htb]
\centering
\includegraphics[width=0.48\textwidth]{cluster_separation.png}
\caption{Bibliometric cluster separation: ML\_DETECTION vs IS\_GOVERNANCE keyword space.}
\label{fig:cluster}
\end{figure}

Bibliometric cluster analysis (keyword-fraction threshold $>$0.6) reveals a structurally
disconnected two-cluster corpus (Table~\ref{tab:clusters}). Keyword co-occurrence confirms
that no ML cluster keyword co-occurs with any governance cluster keyword above count~2 ---
direct bibliometric evidence for the Operationalization Chasm. The three BRIDGING papers
\cite{p008,p013,p043} represent the only cross-citing relations in the corpus; their
existence confirms the chasm is not absolute but also that three papers cannot constitute
a research tradition.

\begin{table}[htb]
\caption{Bibliometric Cluster Structure (N\,=\,45)}
\label{tab:clusters}
\begin{tabular*}{\hsize}{@{\extracolsep{\fill}}lrp{4.2cm}@{}}
\toprule
\textbf{Cluster} & \textbf{N} & \textbf{Characterization} \\
\colrule
ML\_DETECTION & 21 & Deep learning, IF, LOF; private/synthetic data; no IS theory \\
IS\_GOVERNANCE & 21 & Dana Desa, agency theory, institutional controls \\
BRIDGING & 3 & Mixed signals — partial overlap of both clusters \\
\botrule
\end{tabular*}
\end{table}

\subsection{Descriptive Themes}

Open coding of 613 code instances across 45 papers produced 10 descriptive themes
(Table~\ref{tab:dt}). The critical cross-cutting finding: \textbf{28/45 papers (62\%)} contain
no IS-theoretical framing (\texttt{IST-NONE}). The distribution across RQ alignment
reveals a structural imbalance: RQ1 is addressed by 22 papers (DT3) on supervised
detection, yet this majority operates under the label-availability assumption that
collapses in village governance. By contrast, DT2 (unsupervised, 18 papers) represents
the minority that is actually applicable --- a methodological inversion that the corpus
has not previously made explicit. The implication is that the frequently-cited
high-performing methods (supervised: DT3) are deployment-inviable for Dana Desa, while
the under-represented methods (unsupervised: DT2) are the viable candidates.

\begin{table*}[htb]
\caption{Descriptive Theme Summary (DT1--DT10)}
\label{tab:dt}
\begin{tabular*}{\hsize}{@{\extracolsep{\fill}}clrrl@{}}
\toprule
\textbf{ID} & \textbf{Theme Label} & \textbf{N} & \textbf{RQ} & \textbf{Representative Papers} \\
\colrule
DT1 & Corruption as computational signals & 16 & 2 & P012, P016, P035, P043, P044 \\
DT2 & Unsupervised ML anomaly detection   & 18 & 1 & P001, P003, P015, P037, P041 \\
DT3 & Supervised detection (private sector) & 22 & 1 & P003, P007, P022, P026, P039 \\
DT4 & Label scarcity / ground truth gap   & 19 & 3 & P001, P003, P015, P022, P037 \\
DT5 & Village fund governance \& corruption & 26 & 2,3 & P012, P013, P016, P035, P044 \\
DT6 & IS-theoretical framing               & 17 & 1,3 & P013, P022, P039, P043, P044 \\
DT7 & Real-time detection gap              & 4  & 3 & P003, P016, P041 \\
DT8 & Developing-country constraints       & 18 & 3 & P008, P012, P013, P016, P035 \\
DT9 & Graph / network-based methods        & 20 & 1 & P003, P014, P015, P022, P026 \\
DT10 & Explainability \& audit requirements & 7 & 1,3 & P003, P016, P039, P041 \\
\botrule
\end{tabular*}
\end{table*}

\subsection{Analytical Themes}

Cross-paper relationship analysis (121 inter-paper relations) produced four analytical
themes. Figure~\ref{fig:keywords} shows the keyword frequency distribution by cluster,
visualising the lexical separation that underlies each theme.

\begin{figure}[htb]
\centering
\includegraphics[width=0.48\textwidth]{keyword_frequency.png}
\caption{Top keyword frequencies by cluster (ML\_DETECTION vs IS\_GOVERNANCE).}
\label{fig:keywords}
\end{figure}

\textbf{AT1 --- The Operationalization Chasm}: The corpus bifurcates into two
non-communicating traditions: 23 papers develop ML methods without governance engagement
(DT2, DT3, DT9); 26 papers document village corruption without computational approaches
(DT5). Nine bridging papers exist but none unite both traditions in a single study design.
Six silencing relations confirm the absence: no paper applies ML to Dana Desa data; no
paper proposes a village fund feature set; no paper maps corruption typology to
ML-detectable signals. The governance implication is severe: the evidence base that
policy-makers rely on to justify ML adoption for public financial integrity cannot actually
deliver deployable tools without a primary study that crosses this divide.

\textbf{AT2 --- The Scalability Illusion}: Technical papers report AUC-ROC $>$0.95
under assumptions that do not hold in village governance: centralised databases (4 papers),
labelled data (12 papers, \texttt{AC-LABEL}), and high transaction volume. Five papers
validate on synthetic data (\texttt{DS-SYNTH}), creating circular validation that cannot
support generalization claims. The implication for procurement policy is concrete:
adopting any of these systems for Dana Desa auditing based on their reported performance
would constitute a category error --- deploying a system whose evaluation conditions
bear no structural resemblance to the target environment.

\textbf{AT3 --- The Absence of IS Theory}: 28/45 papers (62\%) contain no IS-theoretical
framework (\texttt{IST-NONE}). IS theory appears exclusively in the governance cluster,
never co-occurring with detection methods. No paper evaluates a detection artefact using
the DeLone-McLean IS Success Model \cite{method_delone2003}, Task-Technology Fit, or any
adoption framework. The implication extends beyond academic completeness: without IS
theory, technical papers have no framework for asking whether their system would be used,
understood, or trusted by the auditors who must act on its outputs.

\textbf{AT4 --- The Ground Truth Paradox}: Five papers report high performance on synthetic
or self-labelled data while acknowledging that real-world fraud labels are unavailable ---
a circular validation pattern that cannot support operational deployment claims. The paradox
is self-sealing: acknowledging the label problem provides epistemological cover for using
substitute validation, while the substitute validation sustains the publication metric
(AUC-ROC) that obscures the deployment gap.

\subsection{DSR Framework Matrix}

Table~\ref{tab:dsr} maps 45 papers to DSR cycles $\times$ context levels.

\begin{table}[htb]
\caption{DSR Cycle $\times$ Context Level Matrix}
\label{tab:dsr}
\begin{tabular*}{\hsize}{@{\extracolsep{\fill}}lrrrrr@{}}
\toprule
\textbf{DSR Cycle} & \textbf{Village} & \textbf{Sub-nat.} & \textbf{Nat.} & \textbf{XN-Dev.} & \textbf{XN-Dvlp.} \\
\colrule
DESIGN    & \textbf{0} & \textbf{0} & 3  & 5  & 2 \\
RELEVANCE & 14         & 0          & 5  & 0  & 0 \\
RIGOR     & 11         & 1          & 7  & 8  & 5 \\
\colrule
\textit{Total} & 14 & 2 & 8 & 8 & 5 \\
\botrule
\end{tabular*}
\end{table}

The \textbf{DESIGN $\times$ village cell contains zero papers} --- the most structurally
significant finding of this review. No study has designed and evaluated an ML artefact
using real village-level government financial data. The implication is not that such
studies are difficult --- it is that the research community has not attempted them.
RELEVANCE$\times$village = 14 confirms that governance scholars have fully established
the problem domain; RIGOR$\times$village = 11 confirms that methodological frameworks
exist. The DESIGN cycle is the missing link, and its absence is robust across all three
sensitivity thresholds.

\subsection{Gap Matrix}

\begin{table*}[htb]
\caption{Integrated Research Gap Matrix}
\label{tab:gaps}
\begin{tabular*}{\hsize}{@{\extracolsep{\fill}}clrp{7cm}@{}}
\toprule
\textbf{ID} & \textbf{Severity} & \textbf{N} & \textbf{Gap Statement} \\
\colrule
G1 & CRITICAL     & 14 & No ML artefact applies to Dana Desa financial data (DESIGN $\times$ village = 0) \\
G2 & CRITICAL     & 19 & Label scarcity eliminates supervised approaches; unsupervised methods remain methodologically underspecified for village governance \\
G3 & CRITICAL     & 40 & No validated feature engineering framework exists for village fund fraud; budget absorption features have $<$5 papers of foundation \\
G4 & PARTIAL      & 37 & IS theory absent from 62\% of detection papers; no D\&M IS Success Model application found \\
G5 & METHODOLOGICAL & 7 & Real-time/streaming deployment architectures absent; explainability insufficient for auditor adoption \\
\botrule
\end{tabular*}
\end{table*}

Gaps G1, G2, and G3 are structurally interconnected in a dependency chain, not merely a
list: G3 (no feature set) is a necessary prerequisite for G1 (no application), because an
ML artefact applied to village data without a theoretically grounded feature set would
produce anomaly scores with no interpretive anchor for auditors. G2 (no labels) determines
method viability: without ground truth, ensemble consensus must substitute for accuracy
metrics, and this methodological consequence must be designed into the primary study from
the outset rather than treated as a limitation. G4 and G5 are second-order but
operationally binding: a technically functional system that lacks IS-theoretical grounding
(G4) or outputs uninterpretable anomaly scores (G5) will not be adopted by APIP units
regardless of detection performance.

\subsection{Sensitivity Analysis}

Results across thresholds T1 (N=45), T2 (N=23), T3 (N=14) confirm that core findings
are robust: the two-cluster gap structure (AT1), IS theory absence (AT3), and
DESIGN$\times$village empty cell remain present regardless of threshold. DT5 (village
governance) shows threshold instability due to village governance papers appearing
predominantly in unranked Indonesian journals, confirming that the domain-override
protocol was methodologically necessary.

% =====================================================================
\section{Discussion}
\label{sec:discussion}
% =====================================================================

\subsection{Answering the Research Questions}

\textbf{RQ1} (ML methods): Isolation Forest \cite{p015}, LOF \cite{p041}, and
Autoencoder-based methods \cite{p003,p015} demonstrate the strongest fit for unlabelled,
low-volume village fund data conditions. Supervised methods dominate the corpus
\cite{p003,p007,p022,p026,p039} but require labels unavailable in the Dana Desa context.
GNNs offer superior relational detection \cite{p014,p045} but presuppose transaction-graph
infrastructure absent from current village records. The practical decision rule emerging
from this synthesis is that method selection for a primary study must be driven by
deployment constraints --- label availability, data volume, interpretability requirements
--- rather than by AUC-ROC rankings in the corpus.

\textbf{RQ2} (Feature engineering): DT1 (16 papers) establishes five transferable
feature categories: single-bidder ratios, amendment frequencies, price escalation, timing
concentration, and vendor diversity \cite{p008,p043}. Two village-specific additions
emerge from DT5: budget absorption anomalies and disbursement timing irregularities
\cite{method_kpk2023}. No prior study integrates these into a unified village fund
detection framework. The implication is that feature engineering for a primary study
cannot be borrowed from the existing corpus --- it must be constructed from corruption
typology documentation and validated against KPK case evidence.

\textbf{RQ3} (Structural gaps): Three CRITICAL gaps preclude direct application
(G1, G2, G3), confirmed through convergent triangulation of thematic synthesis,
bibliometric disconnection, and DSR empty cells. The Operationalization Chasm is not
merely a practical gap --- it reflects an epistemological difference: for the ML
tradition, corruption is a statistical anomaly in a feature space; for the governance
tradition, it is a social act within a principal-agent structure
\cite{method_jensen1976}. Bridging these requires an explicit translation layer --- a
corruption typology mapped to computable features --- that neither tradition has produced.

\subsection{Theoretical Contributions}

AT3's finding --- 62\% IS theory absence --- carries an implication that extends beyond
this corpus. The ML fraud detection field has optimised for AUC-ROC on benchmark
datasets, a metric orthogonal to the IS success dimensions (information quality, system
quality, net benefits \cite{method_delone2003}) that determine whether a detection system
can function as a governance instrument. This IS--ML gap is not a matter of academic
completeness --- it is the reason why technically performant systems remain undeployed:
without IS-theoretical grounding, there is no framework for designing the human-system
interface, the institutional workflow, or the explanation format that auditors require to
act on anomaly outputs. This SLR thus motivates a DSR-positioned primary study that
evaluates the detection artefact not only by anomaly score distribution but also by
governance utility criteria.

AT1 constitutes a theoretical contribution beyond gap identification: the chasm between
traditions reflects a deeper conceptual problem. Neither tradition has asked what corruption
looks like as a computational feature in village fund governance. The ML tradition lacks
the sociological vocabulary; the governance tradition lacks the algorithmic vocabulary.
Bridging this requires a corruption typology that is simultaneously sociologically valid
--- grounded in KPK modus operandi documentation --- and computationally actionable ---
mappable to observable expenditure features in Siskeudes records.

\subsection{Practical Implications}

The synthesis specifies five concrete design requirements for a Dana Desa detection tool:
(1) unsupervised ensemble architecture (IF + LOF + Autoencoder) to address G2;
(2) feature engineering grounded in the seven-typology framework derived from DT1 + DT5
to address G3; (3) annual batch processing of village disbursement records to match
current Siskeudes reporting cycles; (4) ranked anomaly output in auditor-interpretable
format to address G5; and (5) IS success evaluation against KPK/BPK surrogate ground
truth to address G4. The top-50 anomaly flagging protocol provides a measurable
effectiveness criterion (search space reduction ratio) without requiring prosecutorial
fraud confirmation --- a pragmatic operationalization of G2's label scarcity constraint.

% =====================================================================
\section{Conclusion}
\label{sec:conclusion}
% =====================================================================

This SLR synthesised 45 papers to map the evidence base for ML-based financial anomaly
detection and identify structural gaps that preclude current application to Dana Desa
village fund governance. Four analytical themes characterise the literature: the
Operationalization Chasm (no study bridges ML detection with village governance data),
the Scalability Illusion (AUC-ROC performance conditioned on inapplicable assumptions),
the Absence of IS Theory (62\% of detection papers contain no IS-theoretical framework),
and the Ground Truth Paradox (circular synthetic-data validation). Five structured gaps
in the integrated gap matrix --- three CRITICAL, one PARTIAL, one METHODOLOGICAL ---
form a dependency chain that motivates a primary empirical study applying an unsupervised
ensemble to real Jambi Province disbursement data (2023--2025).

The DSR framework matrix delivers the most structurally significant finding:
DESIGN$\times$village = 0 across all 45 papers and all three sensitivity thresholds.
This zero cell is not a minor gap in a mature literature --- it is an undiscovered
continent. The RELEVANCE knowledge exists (14 papers); the RIGOR methods exist
(11 papers). What has never been attempted is the artefact that connects them.

The primary theoretical contribution is characterising the nature of the divide between
ML detection and IS governance traditions as epistemological, not merely empirical.
Resolving the Operationalization Chasm requires constructing a theoretically grounded
bridge between governance corruption typology and computational feature design --- a bridge
whose absence this SLR documents with precision and whose construction the primary study
delivers.
construction the primary study delivers.

\textbf{Limitations}: absence of pre-registered protocol; potential underrepresentation of
Chinese-language GNN literature; DT5 quality scoring disadvantage for Indonesian journals.
\textbf{Future research}: (1) apply the unsupervised ensemble to Jambi Province Dana Desa
data; (2) develop IS Success Model evaluation criteria for detection artefact assessment;
(3) extend toward real-time monitoring integration with SiKPA/SIMDA village financial
systems; (4) replicate across provinces and comparable international fiscal
decentralisation programmes.

%% ─────────────────────────────────────────────────────────────────────────────
\section*{Acknowledgements}
Authors would like to thank the reviewers for their constructive feedback. [Full
acknowledgements to be restored in camera-ready version.]

%% ─────────────────────────────────────────────────────────────────────────────
\bibliographystyle{elsarticle-num}
\bibliography{references}

\end{document}
